% ============================================================
%   Chapter 4  \textemdash\  System Design
% ============================================================

\chapter{System Design}\label{ch:design}

This chapter describes how the concepts introduced in Chapter~\ref{ch:sota} are instantiated into
a concrete software architecture. Section~\ref{sec:design-architecture} presents the layered
wrapper chain, Sections~\ref{sec:design-obs}--\ref{sec:design-action} formalise the observation
and action spaces, Section~\ref{sec:design-agent} describes the PPO agent,
Section~\ref{sec:design-shields} the two safety shields,
Section~\ref{sec:design-reward} the reward shaper, Section~\ref{sec:design-curriculum} the
curriculum manager and Section~\ref{sec:design-metrics} the live metrics pipeline.

\section{Architectural overview}\label{sec:design-architecture}

The framework is organised as a vertical stack of components that communicate through the
standard \texttt{gymnasium.Env} interface. Each layer consumes the observation produced by the
layer below and (possibly) transforms both the action travelling downwards and the reward
travelling upwards. Figure~\ref{fig:wrapper-chain} shows the chain.

\begin{figure}[h]
  \centering
  \small
  \begin{tabular}{c}
    \hline
    \textbf{PPOAgent}                            \\[2pt]
    \(\downarrow\) proposed action $\; a_t$ \qquad \(\uparrow\) observation $\; s_t$, reward $\; r_t$ \\[4pt]
    \textbf{Safety Shield}\; (\texttt{none}\,/\,\texttt{basic}\,/\,\texttt{adaptive}) \\[2pt]
    \(\downarrow\) executed action $\; a_t^\star$ \qquad \(\uparrow\) observation, base reward \\[4pt]
    \textbf{CarlaRewardShaper}                  \\[2pt]
    \(\downarrow\) executed action $\; a_t^\star$ \qquad \(\uparrow\) shaped reward $\; \tilde{r}_t$ \\[4pt]
    \textbf{CarlaEnv}\; (\texttt{gymnasium.Env} wrapping the CARLA Python API) \\[2pt]
    \(\downarrow\) vehicle command \qquad \(\uparrow\) sensor readings \\[4pt]
    \textbf{CARLA Simulator}\; (\texttt{CarlaUE4.exe}) \\[2pt]
    \hline
  \end{tabular}
  \caption{Wrapper chain of the framework. Each layer is a \texttt{gymnasium.Wrapper} that can be
           composed, replaced or removed without touching its neighbours.}
  \label{fig:wrapper-chain}
\end{figure}

A subtlety of this layout deserves an explicit comment. The safety shield is placed \emph{above}
the reward shaper, not below it. This ordering has a functional reason: the shaper reads the
\texttt{info} keys \texttt{executed\_action} and \texttt{proposed\_action}\textemdash both
produced by the shield\textemdash to compute components such as the shield-intervention penalty,
the smoothness penalty during an override and the lane-change transition guard. If the shaper
were placed above the shield, those keys would not yet exist when the shaper inspects the
\texttt{info} dictionary, and the shield interaction could not be accounted for in the reward.

The environment runs in \emph{synchronous mode} with a fixed time step of \(\Delta t =
0.05\,\mathrm{s}\). Every call to \texttt{env.step(action)} sends the action to CARLA, requests a
single \texttt{world.tick()} and collects the resulting sensor readings. This guarantees
determinism given a fixed seed, at the cost of not being able to use CARLA's asynchronous
rendering speed-ups.

\section{Observation space}\label{sec:design-obs}

The observation vector \(s_t \in \mathbb{R}^{735}\) concatenates three semantic LIDAR channels and
a small block of geometric and route features. Appendix~\ref{app:observation} contains the full
index-by-index table; this section describes the design rationale.

\subsection{Semantic LIDAR}\label{sec:design-obs-lidar}

The framework uses a simulated LIDAR with \(N = 240\) rays uniformly distributed over
\(360^\circ\) (angular resolution \(1.5^\circ/\text{ray}\), range \(50\,\mathrm{m}\)). Each ray
is normalised to the \([0, 1]\) interval, where \(1\) denotes a free ray up to the maximum range
and values close to \(0\) denote an obstacle next to the vehicle. The semantic channel of each
ray is obtained from CARLA's \texttt{SemanticLidarSensor}, which tags every hit with the
\texttt{carla.CityObjectLabel} of the struck primitive.

Three derived channels are packed into the observation:

\begin{itemize}
  \item \(\text{obs}[0:240]\) \textemdash\ \textbf{combined channel}. Includes every obstacle
        (dynamic and static) plus road edges.
  \item \(\text{obs}[240:480]\) \textemdash\ \textbf{dynamic-only channel}. Keeps only rays that
        struck a vehicle or a pedestrian.
  \item \(\text{obs}[480:720]\) \textemdash\ \textbf{static-only channel}. Keeps only rays that
        struck a wall, a barrier, a pole or a similar static object.
\end{itemize}

The decomposition is motivated by an empirical observation: a single-channel LIDAR induces a
\emph{zig-zag} pattern when the vehicle drives along a guardrail, because the policy interprets
the low returns from the guardrail as a moving obstacle and keeps trying to evade it laterally.
Separating dynamic from static obstacles allows the policy to learn that static obstacles can be
treated as soft boundaries while dynamic obstacles require anticipation.

\subsection{Lane, vehicle and route features}\label{sec:design-obs-vector}

A \(15\)-dimensional vector block is appended to the LIDAR. Its components are:

\begin{itemize}
  \item \(\text{obs}[720:728]\) \textemdash\ \textbf{lane features (8 dims)}:
        \(\texttt{lateral\_offset\_norm}\),
        \(\texttt{heading\_error\_norm}\),
        \(\texttt{on\_edge\_warning}\),
        \(\texttt{lane\_width\_norm}\),
        \(\texttt{dist\_left\_edge\_norm}\),
        \(\texttt{dist\_right\_edge\_norm}\),
        \(\texttt{lane\_change\_permitted}\),
        \(\texttt{road\_curvature\_norm}\).
  \item \(\text{obs}[728:730]\) \textemdash\ \textbf{vehicle state (2 dims)}:
        \(\texttt{speed\_norm}\), \(\texttt{steering}\).
  \item \(\text{obs}[730:735]\) \textemdash\ \textbf{route features (5 dims)}:
        \(\texttt{next\_wp\_angle\_norm}\) (\(5\,\mathrm{m}\) ahead),
        \(\texttt{wp\_angle\_20m\_norm}\),
        \(\texttt{progress\_norm}\),
        \(\texttt{speed\_limit\_norm}\),
        \(\texttt{speed\_ratio}\).
\end{itemize}

All features are extracted from CARLA's Waypoint API and therefore carry centimetre-level
accuracy. The two offsets \(\texttt{dist\_left\_edge\_norm}\) and
\(\texttt{dist\_right\_edge\_norm}\) are normalised by the lane width, so that both take the
value \(0.5\) when the vehicle is at the centre of the lane and approach \(0\) as it nears the
corresponding edge.

\section{Action space}\label{sec:design-action}

The action space is the continuous box \(\mathcal{A} = [-1, 1]^2\). Its two components are
\begin{itemize}
  \item \(a^{(1)} \in [-1, 1]\) \textemdash\ \emph{steering}, mapped linearly to
        \(\delta \in [-\delta_{\max}, +\delta_{\max}]\) with \(\delta_{\max} = 0.6\,\mathrm{rad}\).
  \item \(a^{(2)} \in [-1, 1]\) \textemdash\ \emph{joint throttle\textendash brake}: positive
        values drive the throttle, negative values drive the brake, zero is coast. The
        acceleration saturation values are \(+3.0\,\mathrm{m/s^2}\) and \(-7.0\,\mathrm{m/s^2}\).
\end{itemize}

Merging throttle and brake into a single signed channel has two practical advantages: it reduces
the dimensionality of the policy output, and it prevents the agent from learning the
counter-productive behaviour of activating throttle and brake simultaneously.

Continuous actions are sampled from a diagonal Gaussian and squashed through a hyperbolic
tangent to keep them inside the box:
\begin{equation}
  u_t \sim \mathcal{N}\!\big(\mu_\theta(s_t),\, \sigma_\theta^2\big),
  \qquad
  a_t = \tanh(u_t).
  \label{eq:policy-squash}
\end{equation}
The log probability of the squashed sample requires the log-determinant of the
tanh Jacobian~\cite{haarnoja2018sac}:
\begin{equation}
  \log \pi_\theta(a_t \mid s_t)
    \;=\; \log \mathcal{N}\!\big(u_t;\, \mu_\theta(s_t),\, \sigma_\theta\big)
        \;-\; \sum_{i} \log\!\big(1 - a_{t,i}^2 + \varepsilon\big),
  \label{eq:logprob-tanh}
\end{equation}
where \(\varepsilon\) is a small constant added for numerical stability. The same correction is
used to evaluate the log probability of an executed action, which is a subtle but important
point for credit assignment under a shield (Section~\ref{sec:impl-shield-credit}).

\section{Agent design}\label{sec:design-agent}

\subsection{Actor\textendash Critic network}\label{sec:design-agent-net}

The actor and the critic share a representation of the observation but have separate heads.
Feeding the raw \(735\)-dimensional observation directly into a shallow MLP dilutes the signal:
the \(720\) LIDAR dimensions dominate the input while the \(15\) lane/vehicle/route features
carry most of the task-relevant information. The network therefore includes a dedicated
\emph{LIDAR encoder}:

\begin{equation}
  \mathrm{enc}(s_t)
    \;=\; \mathrm{Tanh}\!\Big( W_2 \, \mathrm{Tanh}\big(W_1\, s_t^{\mathrm{LIDAR}} + b_1\big) + b_2 \Big)
    \;\in\; \mathbb{R}^{64},
  \label{eq:lidar-encoder}
\end{equation}
with \(W_1 \in \mathbb{R}^{256\times 720}\) and \(W_2 \in \mathbb{R}^{64 \times 256}\). The
encoder output is concatenated with the \(15\)-dimensional vector block to produce a
\(79\)-dimensional feature \(\phi(s_t)\):

\begin{equation}
  \phi(s_t) \;=\; \big[\, \mathrm{enc}(s_t)\; ;\; s_t^{\mathrm{vector}} \,\big] \in \mathbb{R}^{79}.
  \label{eq:features}
\end{equation}

The critic is a two-layer MLP with \(256\) units per hidden layer and a scalar output:
\begin{equation}
  V_\theta(s_t) \;=\; W^{V}\, \mathrm{Tanh}\!\Big( \mathrm{MLP}_{\mathrm{critic}}\big(\phi(s_t)\big) \Big) + b^{V}.
  \label{eq:critic-head}
\end{equation}

The actor shares the same trunk layout but, instead of a scalar output, produces a
two-dimensional mean vector and has a separate log-standard-deviation parameter that does not
depend on the state:
\begin{equation}
  \mu_\theta(s_t) \;=\; W^{\mu}\, \mathrm{Tanh}\!\Big( \mathrm{MLP}_{\mathrm{actor}}\big(\phi(s_t)\big) \Big) + b^{\mu},
  \qquad
  \log \sigma_\theta \in \mathbb{R}^2.
  \label{eq:actor-head}
\end{equation}

The log-standard-deviation is clamped to the range \([-3.0,\, -0.2]\). Its initial value is set to
\(-0.7\), which corresponds to a pre-\(\tanh\) standard deviation of approximately \(0.50\). This
gives broad but non-saturating exploration at the beginning of training.

\subsection{PPO update}\label{sec:design-agent-ppo}

At each environment step the agent appends the tuple
\(\big(s_t,\, a^\star_t,\, \log \pi_\theta(a^\star_t\mid s_t),\, r_t,\, d_t,\, \mathrm{trunc}_t,\, V^{\text{boot}}_t\big)\)
to the rollout buffer. Two details are specific to this project:

\begin{enumerate}
  \item The buffer stores the \emph{executed} action \(a^\star_t\) (the action that actually
        reached the simulator), not the action sampled from the policy. The log probability
        stored in the buffer is recomputed by \texttt{agent.evaluate\_action} with
        Equation~\eqref{eq:logprob-tanh}. This is what ties the gradient computation to the
        action the environment saw, avoiding the credit-assignment error that would otherwise be
        introduced by the shield.
  \item When an episode is truncated by a timeout but not terminated (no crash, no off-road),
        the value \(V^{\text{boot}}_t = V_\theta(s_{t+1})\) is stored and used to bootstrap GAE
        at that step. When the episode terminates naturally, \(V^{\text{boot}}_t = 0\) is stored.
\end{enumerate}

When the buffer reaches \(M = 2048\) samples, GAE is run backwards through the rollout
(Equation~\eqref{eq:gae}), advantages are standardised to zero mean and unit variance, and the
return target is divided by the running standard deviation described in
Section~\ref{sec:design-returnnorm}. A loop of up to \(K = 10\) epochs of PPO is then performed,
with an early stop triggered by Equation~\eqref{eq:kl-schulman} once the moving average of the
approximate KL divergence exceeds \(D^{\star}_{\mathrm{KL}}\). The gradient is clipped to an
\(\ell_2\) norm of \(1.0\) and the learning rate follows the cosine schedule of
Equation~\eqref{eq:cosine-lr}.

\subsection{Observation normalisation}\label{sec:design-obsnorm}

Observation normalisation is performed online with Welford's incremental mean/variance
estimator~\cite{welford1962notes}. Given a running mean \(\hat{\mu}\), running variance
\(\hat{\sigma}^2\) and count \(N\), the update for a batch \(\mathbf{x}\) of size \(B\) is
\begin{align}
  \Delta \;&=\; \bar{\mathbf{x}} - \hat{\mu}, \\
  \hat{\mu} \;&\leftarrow\; \hat{\mu} + \Delta\, \frac{B}{N+B}, \\
  \hat{\sigma}^2 \;&\leftarrow\; \frac{N\hat{\sigma}^2 + B \, \mathrm{Var}(\mathbf{x}) + \Delta^2\, \frac{NB}{N+B}}{N+B}, \\
  N \;&\leftarrow\; N + B.
  \label{eq:welford}
\end{align}
Normalised observations are clipped to \([-5, 5]\) to contain outliers.

An implementation detail that turned out to be important is that normalisation is applied
\emph{only} to the \(15\)-dimensional vector block, not to the \(720\) LIDAR dimensions. The
LIDAR is already in \([0,1]\) by construction, so normalisation would not change its scale but
would artificially introduce outliers when the curriculum advances from stage 0 (no NPCs, the
dynamic channel is the constant \(1.0\)) to stage 1 (dynamic channel suddenly has non-unit
entries). Keeping the LIDAR un-normalised avoids this distribution shift and therefore avoids
the gradient spikes observed in early prototypes.

\subsection{Return normalisation}\label{sec:design-returnnorm}

The return targets passed to the critic loss in Equation~\eqref{eq:ppo-aux} can span several
orders of magnitude, which makes MSE training unstable. To counter this, the returns are divided
by the running standard deviation
\begin{equation}
  \hat{R}_t \;\leftarrow\; \frac{\hat{R}_t}{\hat{\sigma}_R + \varepsilon},
  \label{eq:return-norm}
\end{equation}
where \(\hat{\sigma}_R\) is computed by the same Welford estimator applied to the rollout returns.
Note that the mean is \emph{not} subtracted; subtracting a running mean of returns would bias the
advantages, which are already standardised separately.

\section{Safety shields}\label{sec:design-shields}

Both shields implement the function \(\Sigma\) of Equation~\eqref{eq:shield-fn}. They differ in
the semantic information they consume and in the strategy by which they generate the corrected
action.

\subsection{Basic shield}\label{sec:design-basic-shield}

The basic shield (\texttt{CarlaSafetyShield}) operates in two layers. The first layer analyses
the LIDAR scan for imminent obstacles; the second layer uses the Waypoint API to detect
unsafe lane-keeping conditions.

\paragraph{Obstacle analysis.}
Let \(\mathrm{scan}_t \in [0,1]^{240}\) be the combined LIDAR channel at time \(t\). The shield
partitions it into three angular sectors:
\begin{align}
  \mathrm{front}_t \;&=\; \mathrm{scan}_t\big[225^\circ\!:\!315^\circ\big] \;\cup\; \mathrm{scan}_t\big[0^\circ\!:\!22.5^\circ\big] \cup \mathrm{scan}_t\big[337.5^\circ\!:\!360^\circ\big], \label{eq:shield-front}\\
  \mathrm{side}^{\mathrm{R}}_t \;&=\; \mathrm{scan}_t\big[60^\circ\!:\!120^\circ\big], \label{eq:shield-rside}\\
  \mathrm{side}^{\mathrm{L}}_t \;&=\; \mathrm{scan}_t\big[240^\circ\!:\!300^\circ\big], \label{eq:shield-lside}
\end{align}
and declares the situation unsafe whenever any of the following conditions holds:
\begin{equation}
  \min \mathrm{front}_t < \tau_{\mathrm{f}}, \quad
  \min \mathrm{side}^{\mathrm{R}}_t < \tau_{\mathrm{s}}, \quad
  \min \mathrm{side}^{\mathrm{L}}_t < \tau_{\mathrm{s}}.
  \label{eq:shield-obstacle}
\end{equation}
The default thresholds are \(\tau_{\mathrm{f}} = 0.15\) and \(\tau_{\mathrm{s}} = 0.04\) in
normalised LIDAR units.

\paragraph{Lane-keeping analysis.}
The shield reads the lateral offset \(\ell_t = \texttt{lateral\_offset\_norm}\) and the heading
error \(h_t = \texttt{heading\_error\_norm}\) from the \texttt{info} dict and declares the
situation unsafe when
\begin{equation}
  |\ell_t| > \tau_\ell
  \quad \text{or} \quad
  |h_t| > \tau_h,
  \label{eq:shield-lane}
\end{equation}
with defaults \(\tau_\ell = 0.82\) (expressed as a fraction of the lane half-width) and
\(\tau_h = 0.60\).

\paragraph{Correction law.}
When the situation is unsafe the shield overrides the action with a proportional controller.
Let \(k_{\mathrm{s}}\) and \(k_{\mathrm{b}}\) be the steering and brake gains. The correction is
\begin{equation}
  \delta^\star_t \;=\; \mathrm{clip}\!\Big(
      -k_{\mathrm{s}}\, \ell_t \;-\; 0.45\, k_{\mathrm{s}}\, h_t \;+\; \mathrm{sidepush}_t,\;
      -1,\, 1\Big),
  \label{eq:shield-steer-correction}
\end{equation}
\begin{equation}
  (tb)^\star_t \;=\; \mathrm{clip}\!\Big(
      -\, k_{\mathrm{b}}\, \max\!\big(0,\; \tau_{\mathrm{f}} - \min\mathrm{front}_t\big) / \tau_{\mathrm{f}},\;
      -1,\, 0\Big),
  \label{eq:shield-brake-correction}
\end{equation}
where \(\mathrm{sidepush}_t\) contains the lateral correction triggered by close obstacles in
\(\mathrm{side}^{\mathrm{R}}\) or \(\mathrm{side}^{\mathrm{L}}\). Equations
\eqref{eq:shield-steer-correction}--\eqref{eq:shield-brake-correction} follow the
\emph{minimal-interference principle}: if \(\ell_t\), \(h_t\) and \(\min\mathrm{front}_t\) are
all far from their thresholds, all three contributions vanish and the shield produces the
zero-override action. Conversely, the closer the system is to an unsafe region, the stronger the
corrective action.

\subsection{Adaptive-horizon shield}\label{sec:design-adaptive-shield}

The adaptive-horizon shield (\texttt{CarlaAdaptiveHorizonShield}) extends the basic design with
three mechanisms: a risk classifier, a trajectory prediction with the bicycle model, and a
priority-ordered candidate search.

\paragraph{Risk classifier.}
Three risk levels are defined: \(\mathrm{safe}\), \(\mathrm{warning}\) and \(\mathrm{critical}\).
Each level fixes a verification horizon \(H\) (number of bicycle-model steps) and a threshold
multiplier \(\kappa\) applied to the basic thresholds. The configuration is
\[
  \mathrm{safe}:\; H=1,\; \kappa = 1.0;
  \qquad
  \mathrm{warning}:\; H=5,\; \kappa = 1.5;
  \qquad
  \mathrm{critical}:\; H=10,\; \kappa = 2.0.
\]
The level is chosen as the worst of two independent criteria. The first is the
\emph{frontal-distance} criterion: let \(d^{\mathrm{dyn}}_t = \min\text{front}^{\mathrm{dynamic}}_t\)
be the minimum of the dynamic-only LIDAR front sector. Then
\begin{equation}
  \mathrm{level}^{\mathrm{frontal}}_t \;=\;
     \begin{cases}
       \mathrm{safe}     & \text{if } d^{\mathrm{dyn}}_t > 0.50,\\
       \mathrm{warning}  & \text{if } 0.20 < d^{\mathrm{dyn}}_t \le 0.50,\\
       \mathrm{critical} & \text{if } d^{\mathrm{dyn}}_t \le 0.20.
     \end{cases}
  \label{eq:shield-risk-front}
\end{equation}
The second criterion, \emph{lateral position}, thresholds \(|\ell_t|\) at \(0.70\) and \(0.85\)
(warning and critical, respectively). If the Waypoint API reports that a lane change is
permitted at the current location, a \emph{critical} lateral level is downgraded to
\emph{warning} so that the shield does not block a legitimate lane-change manoeuvre.

\paragraph{Trajectory check.}
Given a candidate action \(a = (\delta, tb)\), the shield asks the bicycle model of
Section~\ref{sec:sota-bicycle} for the trajectory \(\{(x_k, y_k, \theta_k)\}_{k=0}^{H}\). Each
predicted position \((x_k, y_k)\) is then projected onto the road using
\texttt{carla.Map.get\_waypoint} restricted to driving lanes. The candidate is declared safe
when all of the following hold:

\begin{equation}
  \min \mathrm{front}^{\mathrm{combined}}_t \;\ge\; \tau_{\mathrm{f}} / \kappa,
  \qquad
  \min \mathrm{side}^{\mathrm{R/L,\,static}}_t \;\ge\; \tau_{\mathrm{s}} / \kappa,
  \label{eq:shield-lidar-check}
\end{equation}
\begin{equation}
  \mathrm{nearest\_pedestrian\_m}_t \;\ge\; 4.0\,\mathrm{m},
  \label{eq:shield-ped-check}
\end{equation}
\begin{equation}
  \forall\, k \in \{1,\ldots,H\}:\;
    |\ell_k| \;\le\; \max\!\big( \tau_\ell / \kappa,\; 0.45 \big),
  \label{eq:shield-lane-check}
\end{equation}
where \(\ell_k = \texttt{lat\_offset}(x_k, y_k) / (\texttt{lane\_width}/2)\) is the normalised
lateral offset of the predicted position with respect to the nearest driving lane. The floor of
\(0.45\) in Equation~\eqref{eq:shield-lane-check} prevents the lateral threshold from becoming
so tight in the \emph{critical} level that no correction action can satisfy it. When the
\emph{lane-change permitted} flag is set, the lateral threshold is further raised to \(1.2\) so
that the transition band between lanes is admissible.

\paragraph{Priority-ordered candidate search.}
If the agent's action \(a_t\) does not pass the trajectory check, the shield inspects an
ordered list of candidate actions and returns the first one that does. The ordering depends on
the dominant threat, which is identified in three mutually exclusive buckets:

\begin{itemize}
  \item \textbf{Pedestrian emergency}\; (\(\mathrm{nearest\_pedestrian\_m}_t < 4\,\mathrm{m}\)).
        The override is always \((\delta, tb) = (0, -1)\) (maximum brake, no search).
  \item \textbf{Dynamic threat}\; (nearby vehicle in front, predicted time-to-contact below
        \(\max(5\,\mathrm{s},\, 1.5\,v_t)\) metres). The first candidates privilege braking;
        steering is adjusted afterwards.
  \item \textbf{Static threat}\; (static obstacle in the side sector). The first candidates
        privilege a steering correction back to the lane centre, accompanied by a mild brake.
  \item \textbf{Lateral-only threat}\; (\(|\ell_t| > 0.65\) without any dynamic or static
        threat). The first candidates apply a lane-centring steering with a small throttle to
        keep enough speed for directional control.
  \item \textbf{Balanced fallback}\; (default): a symmetric list of brake/steer combinations.
\end{itemize}

If none of the candidates pass the trajectory check, the shield returns the \emph{fallback}
action
\begin{equation}
  a_t^{\mathrm{fallback}} \;=\; \big(\mathrm{clip}(-1.2\,\ell_t,\, -0.8,\, 0.8),\; -0.4\big).
  \label{eq:shield-fallback}
\end{equation}

\subsection{Bicycle-model parameters}\label{sec:design-bicycle-params}

The bicycle model described in Section~\ref{sec:sota-bicycle} is instantiated with the
parameters of CARLA's Tesla Model 3, summarised in Table~\ref{tab:bicycle-params}.

\begin{table}[h]
  \centering
  \caption{Parameters of the kinematic bicycle model used by the adaptive shield.}
  \label{tab:bicycle-params}
  \small
  \begin{tabular}{@{}lll@{}}
    \toprule
    \textbf{Symbol} & \textbf{Value} & \textbf{Meaning} \\
    \midrule
    \(L\)               & \(2.87\,\mathrm{m}\)          & Wheelbase of the Tesla Model 3 \\
    \(\delta_{\max}\)   & \(0.60\,\mathrm{rad}\) (\({\approx}34^\circ\)) & Maximum steering angle \\
    \(\Delta t\)        & \(0.05\,\mathrm{s}\)          & Time step, matched to the simulator \\
    \(a_{\max}\)        & \(3.0\,\mathrm{m/s^2}\)       & Maximum forward acceleration \\
    \(-a_{\min}\)       & \(7.0\,\mathrm{m/s^2}\)       & Maximum deceleration (brake) \\
    \bottomrule
  \end{tabular}
\end{table}

\section{Reward shaping}\label{sec:design-reward}

The reward shaper \texttt{CarlaRewardShaper} produces a shaped reward \(\tilde{r}_t\) by
augmenting the sparse base reward \(r^{\mathrm{base}}_t\) of \texttt{CarlaEnv} (success
\(+30\), collision \(-10\), off-road \(-20\)) with twelve dense components. The shaped reward is
\begin{equation}
  \tilde{r}_t \;=\; r^{\mathrm{base}}_t + r^{\mathrm{alive}}_t + r^{\mathrm{speed}}_t + r^{\mathrm{lane}}_t + r^{\mathrm{head}}_t + r^{\mathrm{prog}}_t + r^{\mathrm{no-shield}}_t
             - p^{\mathrm{smooth}}_t - p^{\mathrm{inv}}_t - p^{\mathrm{edge}}_t - p^{\mathrm{wrong-h}}_t - p^{\mathrm{shield}}_t - p^{\mathrm{idle}}_t - p^{\mathrm{drift}}_t.
  \label{eq:shaped-reward}
\end{equation}

Each component is described below. Unless stated otherwise, all variables read from the
\texttt{info} dictionary come from the Waypoint API.

\paragraph{Speed gate.}
A common factor \(g_t \in [0,1]\) modulates the components that only make sense when the vehicle
is actually moving. It is a linear ramp of the speed \(v_t\) between the two
thresholds \(v_{\min}\) and \(v_g\):
\begin{equation}
  g_t \;=\; \mathrm{clip}\!\left( \frac{v_t - v_{\min}}{v_g - v_{\min}},\, 0,\, 1 \right).
  \label{eq:speed-gate}
\end{equation}
With defaults \(v_{\min} = 5\,\mathrm{km/h}\) and \(v_g = 10\,\mathrm{km/h}\).

\paragraph{Alive bonus.}
\begin{equation}
  r^{\mathrm{alive}}_t \;=\; c_{\mathrm{alive}}\, g_t \;\cdot\; \mathbb{1}[\,\text{on\_road}_t\,].
  \label{eq:r-alive}
\end{equation}
Gating by \(g_t\) is essential to prevent the local optimum \emph{``stop and collect
0.15/step''} that plagued earlier prototypes of the shaper.

\paragraph{Speed bonus.}
Let \(v^{\mathrm{lim}}_t\) be the effective speed limit (from CARLA's speed limits or the
project-level target, whichever is defined). A curvature correction reduces the target speed in
sharp bends:
\begin{equation}
  v^{\mathrm{target}}_t \;=\; v^{\mathrm{lim}}_t \cdot \max\!\Big(\,1 - c_{\mathrm{curv}}\cdot \min\!\big(|\kappa_t|/0.6,\, 1\big),\; 0.4\,\Big).
  \label{eq:target-speed}
\end{equation}
The speed bonus is a Gaussian centred on this curvature-adjusted target, with a speed-ceiling
penalty when the vehicle exceeds the raw limit by more than the margin:
\begin{equation}
  r^{\mathrm{speed}}_t \;=\; w_{\mathrm{v}}\, \exp\!\left( -\frac{(v_t - v^{\mathrm{target}}_t)^2}{2\sigma^2} \right)
     \;-\; w_{\mathrm{v}} \cdot 0.8\, \max\!\Big(\,0,\, \frac{v_t - (1+\eta)\,v^{\mathrm{lim}}_t}{v^{\mathrm{lim}}_t}\,\Big),
  \label{eq:r-speed}
\end{equation}
with \(\sigma = 0.35\, v^{\mathrm{target}}_t\) and \(\eta = 0.05\) the tolerance margin.

\paragraph{Lane centring and heading.}
\begin{equation}
  r^{\mathrm{lane}}_t \;=\; w_{\mathrm{lc}}\, \max(g_t,\, 0.3)\, \mathrm{clip}\!\big(\min(d^{\mathrm{L}}_t, d^{\mathrm{R}}_t)/0.5,\, 0,\, 1\big),
  \label{eq:r-lane}
\end{equation}
\begin{equation}
  r^{\mathrm{head}}_t \;=\; w_{\mathrm{h}}\, g_t \, \exp\!\left( -\frac{h_t^2}{2\cdot 0.40^2} \right),
  \label{eq:r-head}
\end{equation}
where \(d^{\mathrm{L}}_t,\, d^{\mathrm{R}}_t\) are the normalised distances to the left and right
lane edges. The floor of \(0.3\) in Equation~\eqref{eq:r-lane} provides a small centring signal
even when the vehicle is stopped, which helps the policy learn to align before accelerating.

\paragraph{Smoothness penalty.}
\begin{equation}
  p^{\mathrm{smooth}}_t \;=\; w_{\mathrm{sm}}\, \big| \delta_t - \delta_{t-1} \big|.
  \label{eq:p-smooth}
\end{equation}

\paragraph{Invasion, edge and wrong-heading penalties.}
The lane-invasion penalty uses the signal of CARLA's lane-invasion sensor; it is weighted by the
invasion severity (a function of \(|\ell_t|\)) and is suppressed when the invasion is
intentional (large steering magnitude) or when a lane change is permitted. The edge penalty is a
quadratic ramp triggered when the closest edge distance drops below \(0.4\); the wrong-heading
penalty triggers when \(|\text{heading\_error\_deg}_t|\) exceeds \(90^\circ\).

\paragraph{Progress bonus.}
A fixed bonus \(w_{\mathrm{prog}}\) is awarded every time the vehicle crosses a distance
milestone of \(\Delta_{\mathrm{prog}} = 25\,\mathrm{m}\). This avoids the pathological behaviour
of a speed bonus that rewards high speed regardless of whether the vehicle is actually moving
along the route.

\paragraph{Shield interaction.}
Let \(a_t\) be the action proposed by the agent and \(a^\star_t\) the action executed by the
environment. Define the shield-divergence measure
\begin{equation}
  \Delta^{\mathrm{shield}}_t \;=\; \frac{\|a^\star_t - a_t\|_2}{2\sqrt{2}}.
  \label{eq:shield-div}
\end{equation}
The normalisation by \(2\sqrt{2}\) maps the divergence to the unit interval (the maximum
distance between two points of \([-1,1]^2\) is \(2\sqrt{2}\)). The shield-interaction components
are
\begin{equation}
  p^{\mathrm{shield}}_t \;=\; \Delta^{\mathrm{shield}}_t\, w_{\mathrm{sh}}\;
    \mathbb{1}[\,\text{shield\_activated}_t\,],
  \qquad
  r^{\mathrm{no-shield}}_t \;=\; w_{\mathrm{no-sh}}\,
    \mathbb{1}[\,\neg \text{shield\_activated}_t\,].
  \label{eq:shield-reward}
\end{equation}
Either \(w_{\mathrm{sh}}\) or \(w_{\mathrm{no-sh}}\) can be zero, which makes it possible to
penalise shield activations, to reward the agent for not needing the shield, or to run in a
neutral configuration.

\paragraph{Idle penalty and shield grace period.}
An idle penalty discourages the agent from stopping. To prevent the penalty from punishing
legitimate shield-induced stalls, a \emph{grace period} of \(G\) steps is opened the first time
the shield activates in an episode. The counter is re-armed \emph{only} when the shield fires
and the grace period has already expired, which prevents an always-active shield from
suppressing the idle penalty forever.

\paragraph{Drift penalty.}
Asymmetric driving (e.g.\ hugging the right edge) is penalised by
\begin{equation}
  p^{\mathrm{drift}}_t \;=\; w_{\mathrm{dr}}\, \max\!\big(0,\, |d^{\mathrm{L}}_t - d^{\mathrm{R}}_t| - 0.3\big)\;
     \mathbb{1}\!\left[\min(d^{\mathrm{L}}_t, d^{\mathrm{R}}_t) < 0.35\right].
  \label{eq:p-drift}
\end{equation}

\paragraph{Lane-change transition guard.}
When the Waypoint API reports \texttt{lane\_change\_permitted} and \(|\ell_t| > 0.5\), the
penalties \(p^{\mathrm{inv}}_t\), \(p^{\mathrm{edge}}_t\) and \(p^{\mathrm{drift}}_t\) are all
set to zero. This keeps the cost of a legitimate lane-change manoeuvre from blocking the
behaviour during training.

The default weights of all components are listed in
Appendix~\ref{app:hyperparameters}, Table~\ref{tab:hp-reward}.

\section{Curriculum manager}\label{sec:design-curriculum}

The curriculum manager divides the interval \([0,\, N_{\max}]\) of NPC vehicles uniformly in
five stages: \((0,\, N_{\max}/4,\, N_{\max}/2,\, 3N_{\max}/4,\, N_{\max})\). For the default
\(N_{\max}=40\) this yields \((0,\,10,\,20,\,30,\,40)\). The manager keeps two counters: the
number of episodes at the current stage \(E_{\mathrm{stage}}\) and a decay-weighted rollback
counter \(C_{\mathrm{rb}}\).

\paragraph{Advancement rule.}
After \(E_{\mathrm{stage}} \ge E^{\star} = 100\) episodes, if the stage-specific condition of
Table~\ref{tab:curriculum-rules} is satisfied, the manager advances one stage and resets both
counters to zero.

\paragraph{Rollback rule.}
After \(E_{\mathrm{stage}} > E^{\star}\), on every episode the counter \(C_{\mathrm{rb}}\)
is incremented by one if the stage is \emph{collapsing} and decremented by one otherwise (but
not below zero). When \(C_{\mathrm{rb}} \ge C^{\star} = 50\), the manager rolls back one stage
and resets all counters. The decay by \(-1\) on non-collapsing episodes prevents a streak of
bad luck from triggering a rollback that is not warranted by a sustained performance drop.

\begin{table}[h]
  \centering
  \caption{Advancement and rollback thresholds per curriculum stage. All rates are 100-episode
           rolling averages.}
  \label{tab:curriculum-rules}
  \small
  \begin{tabular}{@{}clll@{}}
    \toprule
    \textbf{Stage} & \textbf{NPCs} & \textbf{Advance if} & \textbf{Rollback if} \\
    \midrule
    0 & \(0\)                 & \(\text{offroad\_rate} < 0.20 \wedge \text{avg\_reward} > 0\)
                              & \textemdash\ \\
    1 & \(N_{\max}/4\)        & \(\text{crash\_rate} < 0.25 \wedge \text{offroad\_rate} < 0.15\)
                              & \(\text{crash\_rate} > 0.45 \vee \text{offroad\_rate} > 0.40\) \\
    2 & \(N_{\max}/2\)        & \(\text{crash\_rate} < 0.15\)
                              & \(\text{crash\_rate} > 0.40\) \\
    3 & \(3N_{\max}/4\)       & \(\text{crash\_rate} < 0.10\)
                              & \(\text{crash\_rate} > 0.35\) \\
    4 & \(N_{\max}\)          & \textemdash\ (terminal)
                              & \(\text{crash\_rate} > 0.30\) \\
    \bottomrule
  \end{tabular}
\end{table}

\paragraph{Non-disruptive NPC update.}
When the stage changes, the manager does \emph{not} rebuild the environment. It updates the
attribute \texttt{CarlaEnv.num\_npc\_vehicles} in place; the change takes effect at the next
\texttt{env.reset()}. This avoids the deadlock observed in earlier prototypes that had to tear
down the wrapper chain every time the stage changed.

\section{Live metrics pipeline}\label{sec:design-metrics}

Training metrics are written to a per-run SQLite database at
\texttt{runs/\textless run\_name\textgreater/metrics.sqlite}. The logger
(\texttt{src/Metrics/live\_metrics.py}) opens the database in \emph{Write-Ahead Logging}
(WAL)~\cite{sqlitewal} mode, which allows concurrent reads and writes: the training process
writes while the Streamlit dashboard reads.

Two tables are used, corresponding to the two time axes of the experiment:

\begin{itemize}
  \item \texttt{metrics\_episode}: one row per episode, indexed by episode number. Columns
        include the reward components, the safety rates, the curriculum stage, and the
        LIDAR-derived minimum distances.
  \item \texttt{metrics\_update}: one row per PPO update, indexed by the cumulative timestep.
        Columns include the policy loss, the value loss, the approximate KL, the number of
        epochs that actually ran, the gradient norm and the current learning rate.
\end{itemize}

The dashboard (\texttt{dataVisualizer.py}) reads both tables every few seconds and renders
interactive Plotly charts. Unknown metric keys are rendered as soft-typed line plots, which
keeps the dashboard forward-compatible with future additions to the logger.
